% Related Work for the Introspective Verification paper.
% Style: plain, direct; no em-dashes; minimal colons/quotes; no anthropomorphic verbs.

\section{Related Work}
\label{sec:related}

\paragraph{Process reward models.}
Outcome verifiers score a full solution and improve selection over a base policy \citep{cobbe2021gsm8k}, and process reward models extend this to the step level, which gives denser feedback and better error attribution \citep{uesato2022solving, lightman2023lets}. Step labels come from human annotation \citep{lightman2023lets}, from Monte Carlo rollouts that estimate whether a step still leads to the correct answer \citep{wang2023mathshepherd}, or from consensus filtering between rollouts and a language-model judge \citep{zhang2025lessons}. ProcessBench standardized step-level evaluation and reports that most classification process reward models trail strong critic models on proof-heavy subsets \citep{zheng2024processbench}, a recent survey organizes the area \citep{zheng2025survey}, and later work extends step-level verification to structured and safety-critical reasoning \citep{pronesti2026verifiable}. These models attach a trained scalar head to the final-layer representation. Our verifier keeps a discriminative head but reads the verdict from an intermediate layer, and combines it with the final layer through a residual.

\paragraph{Generative and code-augmented verification.}
The most accurate recent verifiers produce the score by generation. A model writes a critique of the step as next-token prediction \citep{zhang2024genrm}, reasons about the step before judging it \citep{xiong2025stepwiser}, and in some systems generates and runs code before the verdict \citep{zhao2025genprm}. This accuracy comes at a cost that grows with the length of the generated critique and with the number of samples used for majority voting \citep{snell2024scaling, zhao2025genprm}. Dyve reduces the average cost by routing easy steps to a fast check and hard steps to slow generation \citep{zhong2025dyve}. Introspection removes the generation and keeps a single forward pass. It is complementary to these verifiers, and a fixed combination improves a strong generative model on every ProcessBench subset.

\paragraph{Label-efficient and shared-backbone reward models.}
Several methods lower the cost of building a process reward model. Implicit process rewards are read off an outcome-trained model through a log-ratio parameterization \citep{yuan2024free}, FreePRM pseudo-labels steps from final-answer correctness \citep{sun2025freeprm}, and ensemble prompting with reverse verification produces domain-agnostic labels \citep{tan2025aurora}. MetaStone-S1 shares the policy backbone and adds a small self-supervised scoring head for step evaluation \citep{wang2025metastone}. These designs reduce label or parameter cost, yet the score is still taken from the final representation or from generation. Our head reads an intermediate layer, and the readout depth is selected automatically rather than fixed.

\paragraph{Reading correctness from hidden states.}
Hidden states carry information about whether a model's own output is correct, which linear probes recover for factual statements \citep{azaria2023internal, kadavath2022language, burns2022discovering, marks2023geometry}, and steering along these directions changes truthfulness \citep{li2023inference}. Answer correctness and hallucination are predictable from internal states, sometimes before the answer is produced \citep{orgad2024llms, liu2024universal, cencerrado2025noanswer}. Closest to our work, hidden-state signals have been turned into rewards for Best-of-N sampling \citep{guo2025mining} and into a lightweight step verifier that probes a frozen backbone \citep{ni2025reprobe}. That line uses frozen representations and targets answer-level truthfulness or Best-of-N selection. We address step-level process verification on ProcessBench, adapt the backbone with LoRA so the representation itself improves on proof domains, combine an intermediate and a final layer, and reach the accuracy of generative process reward models at a fraction of the compute.

\paragraph{Layer-wise representations, probing, and adaptation.}
Linear classifiers on hidden states are a standard tool for locating information across depth \citep{alain2016understanding}, and lenses map intermediate states to output distributions \citep{belrose2023eliciting}. This work reports that task-relevant structure is often most decodable in intermediate layers rather than the final layer, which motivates our readout depth. Our verdict head is trained with a parameter-efficient adaptation of the backbone \citep{hu2021lora}, and the chosen depth is produced by the model's own measurement rather than set by hand. Test-time verification and self-consistency use such verifiers to select among sampled solutions \citep{wang2022selfconsistency}.
